\documentclass[11pt]{letter}
\usepackage[top=0.75in,bottom=0.75in,left=0.95in,right=0.95in]{geometry}
\usepackage{url}

% The letter class uses a generous \parskip and reserves 6\medskipamount above
% the signature. Both are trimmed so the letter fits on a single page.
\setlength{\parskip}{0.35\baselineskip}
\makeatletter
\renewcommand{\closing}[1]{\par\nobreak\vspace{\parskip}%
  \stopbreaks
  \noindent
  \ifx\@empty\fromaddress\else
  \hspace*{\longindentation}\fi
  \parbox{\indentedwidth}{\raggedright
       \ignorespaces #1\\[1.5\medskipamount]%
       \ifx\@empty\fromsig
           \fromname
       \else \fromsig \fi\strut}%
   \par}
\makeatother

\signature{Nitin Kumar Singh\\
Discipline of Chemical Engineering\\
Indian Institute of Technology (IIT) Gandhinagar\\
Palaj, Gujarat 382355, India\\
singh\_nitin@iitgn.ac.in}

\date{\today}

\pagestyle{empty}

\begin{document}

\begin{letter}{Editors-in-Chief\\
Biophysical Chemistry}

\opening{Dear Editors,}

I am pleased to submit our manuscript titled \textbf{``LigPlot3D: An Interactive Web-Based Framework for Geometric Analysis of Protein--Ligand Interactions''} for consideration as an original research article in \textit{Biophysical Chemistry}.

Characterizing non-covalent interactions at protein--ligand interfaces is fundamental to understanding the physical-chemical principles of molecular recognition, binding affinity, and structural specificity. While molecular docking generates candidate poses, systematically dissecting contact geometries across multiple poses remains a practical bottleneck. Existing interaction profilers are designed primarily for batch workflows, 2D schematics obscure three-dimensional spatial arrangements, and most tools require complex local installations or remote server uploads. This leaves a notable gap for researchers seeking rapid, privacy-preserving, and interactive biophysical evaluation of binding modes directly within the browser.

LigPlot3D addresses this need as an open-source, client-side web application where structure parsing, interaction detection, and 3D molecular rendering execute entirely in the user's browser. Because computations run locally, structural data never leave the client machine, eliminating server dependencies and safeguarding proprietary compounds. LigPlot3D deterministically identifies key non-covalent interactions governing biomolecular complexes (hydrogen bonds, salt bridges, hydrophobic contacts, $\pi$-stacking, cation--$\pi$ interactions, and halogen bonds), presenting them simultaneously through synchronized 3D WebGL scenes and exportable interaction tables.

Several aspects of this work should be of strong interest to the readership of \textit{Biophysical Chemistry}. First, rather than relying on heuristic scoring, each interaction criterion is derived from published geometric conventions with explicit distance and angle cutoffs documented in the interface, ensuring transparent and reproducible contact assignments. Second, the manuscript presents a quantitative benchmark comparing these thresholds against an established interaction profiler (PLIP), discussing how boundary formulations influence contact identification in macromolecular interfaces. Finally, the framework is validated on a benchmark complex with experimentally characterized binding determinants, accurately capturing the aromatic cage, salt bridge, and hydrogen-bonding networks that stabilize the ligand.

This work directly aligns with the scope of \textit{Biophysical Chemistry} by providing a reproducible, accessible method to analyze protein--ligand binding interactions directly in the web browser without requiring any local installation or remote servers. LigPlot3D is freely accessible at \url{https://singhnitink.github.io/LigPlot3D/}, with all source code openly available.

This manuscript is original, has not been published elsewhere, and is not under consideration by any other journal. There are no conflicts of interest to declare. Thank you very much for your time and consideration of this submission.

\closing{Sincerely,}

\end{letter}
\end{document}
